% ============================================================
% 2. MARCO TEÓRICO Y ESTADO DEL ARTE
% Lo que el lector necesita para entender los métodos y los
% resultados, y dónde está la literatura. Solo teoría ajena:
% lo que se desarrolló en el trabajo (p. ej. la distancia
% efectiva de la tarea 24) va en el capítulo 3.
% ============================================================
\chapter{Marco teórico y estado del arte}
\label{cap:marco}

% ------------------------------------------------------------
\section{Difracción escalar}
\label{sec:difraccion}

\pendiente{integral de Rayleigh--Sommerfeld y espectro angular como la misma
solución en dos dominios \cite{goodman2005, ersoy2006}; función de
transferencia $H(f_x,f_y)$; componentes evanescentes.}

% ------------------------------------------------------------
\section{Espectro angular por FFT y sus límites de muestreo}
\label{sec:asm-fft}

\pendiente{discretización; paso en frecuencia $1/L$ fijado por la ventana;
condición de muestreo de la función de transferencia y la frecuencia máxima
que la cumple (Ec.~12 de \cite{zhao2020}); aliasing; la ventana de salida
hereda la de entrada \cite{li2007}.}

% ------------------------------------------------------------
\section{Soluciones conocidas al submuestreo}
\label{sec:soluciones}

\subsection{Espectro angular de banda limitada (BLAS)}
\pendiente{recorte de la función de transferencia \cite{matsushima2009}; por
qué pierde detalle a distancias largas.}

\subsection{Ventana ampliada y muestreo no uniforme}
\pendiente{WWASM \cite{yu2012} y NUFFT \cite{kim2014}, en un párrafo cada uno:
qué resuelven y qué cuestan.}

% ------------------------------------------------------------
\section{Espectro angular por producto matricial (MPASM)}
\label{sec:mpasm}

\pendiente{la DFT como triple producto de matrices; los coeficientes $K_f$ y
$s$ y el muestreo de salida $r$; ancho de banda efectivo frente a BLAS;
complejidad \cite{zhao2020, zhao2020code}.}

\subsection{Relación con la transformada Z \emph{chirp}}
\pendiente{la CZT como otro cálculo de la DFT con muestreo arbitrario y el
parecido estructural con el MPASM (objetivo 2).}

% ------------------------------------------------------------
\section{Microscopía holográfica digital sin lentes}
\label{sec:dlhm}

\pendiente{geometría fuente puntual--muestra--sensor ($z$, $L$); magnificación
$M = L/z$; apertura numérica y resolución; modelo de formación del holograma
como difracción de un frente esférico más magnificación no homogénea
\cite{lopera2024}.}

\subsection{Reconstrucción en DLHM}
\pendiente{RS no aproximada como referencia de exactitud \cite{buitrago2019};
el espectro angular con iluminación de onda esférica que se usa de rutina en
el laboratorio, que es la línea base del objetivo 3.}

% ------------------------------------------------------------
\section{Antecedentes}
\label{sec:antecedentes}

% Base: sección B del anteproyecto.
Las limitaciones de muestreo del espectro angular basado en FFT se
identificaron tempranamente \cite{li2007}. El método de banda limitada
\cite{matsushima2009} restringe el cálculo a las frecuencias que sí cumplen la
condición de muestreo y es hoy la solución más difundida; una alternativa
opuesta amplía la ventana de entrada para comprimir el muestreo en frecuencia
\cite{yu2012}; también se han explorado transformadas de Fourier con muestreo
no uniforme, que amplían el rango útil a costa de una implementación más
compleja \cite{kim2014}.

El MPASM se propuso en 2020 junto con una implementación en Python de acceso
abierto \cite{zhao2020, zhao2020code}. Los resultados publicados lo comparan
contra el ASM convencional, el método de banda limitada y la integral de
Rayleigh--Sommerfeld, cuya evaluación directa sirve de referencia porque no
impone restricciones de muestreo sobre el plano de salida. Toda esa
validación, sin embargo, se hizo en simulación ---un haz gaussiano atravesando
una lente---, sin datos experimentales.

En DLHM, el grupo de Óptica de la Universidad Nacional de Colombia (Sede
Medellín) ha abordado la reconstrucción mediante la implementación no
aproximada de la integral de Rayleigh--Sommerfeld (RSD), que elimina las
restricciones sobre el rango de propagación a costa de un elevado costo
computacional, mitigado en GPU \cite{buitrago2019}; por esa ausencia de
restricciones de muestreo se emplea aquí como referencia de exactitud
numérica. La reconstrucción de rutina en el montaje DLHM del laboratorio, en
cambio, se realiza con un propagador de espectro angular que incorpora la
iluminación de onda esférica del \emph{pinhole}, y es esa la línea base frente
a la cual se contrastan los métodos implementados. Este trabajo se inserta en
esa línea, evaluando una familia de propagadores ---la del producto
matricial--- no explorada antes en DLHM.
